\documentclass[10pt,a4paper]{amsart}
\usepackage[margin=19mm]{geometry}
\usepackage{amsmath,amssymb,mathtools}
\usepackage[scr=rsfs,cal=euler]{mathalfa}
\usepackage{xfrac}
\usepackage[unicode,hidelinks,bookmarks=false]{hyperref}
\newcommand{\ie}{i.e.\ }
\newcommand{\cf}{cf.\ }
\newcommand{\resp}{respectively\ }
\newcommand{\Subsection}{\subsection}
\providecommand{\nobreakdash}{\nobreak\hbox{-}}
\setlength{\emergencystretch}{2em}
\multlinegap0pt
\usepackage{amssymb}
\usepackage{mathtools}
\usepackage[shortlabels]{enumitem}
\usepackage{tcolorbox}
\usepackage{float}
\usepackage{xcolor}
\usepackage{tensor}
\newtheorem{lemma}{Lemma}
\newcommand{\Arg}{\operatorname{Arg}}
\title{Algebraic representatives of the ratios $\zeta(2n+1)/\pi^{2n}$ and $\beta(2n)/\pi^{2n-1}$}
\author{Luc Rams\`es Talla Waffo}
\date{Source: arXiv:2602.16761v4 (CC BY 4.0)}
\begin{document}
\maketitle

\noindent\emph{Source and licence.} Algebraic representatives of the ratios $\zeta(2n+1)/\pi^{2n}$ and $\beta(2n)/\pi^{2n-1}$. Version arXiv:2602.16761v4 (CC BY 4.0), distributed by arXiv under the Creative Commons Attribution 4.0 International licence, http://creativecommons.org/licenses/by/4.0/, as recorded in the arXiv licence metadata for this exact version and shown on the abstract page https://arxiv.org/abs/2602.16761v4. Full licence notice: \texttt{licenses/arxiv-2602.16761-CC-BY-4.0.txt}.\par\smallskip

\noindent\textit{Editorial scope.} Counted unit: the article's lemma lemma (source0.tex lines 470--487). The article's complete original proof of it is delivered with it (source0.tex lines 489--680, one \texttt{proof} environment of 192 lines). No mathematics is written, repaired or re-derived: the statement and the proof are the article's own lines, in the article's own notation, and no retained line is substituted or re-flowed.  No further source-local context is needed: every cross-reference of every retained line resolves inside the two delivered spans.  Nothing was omitted from any retained span, so the package carries no omission record.  No retained line contains a citation, so the wrapper carries no bibliography and no bibliography entry.  No retained line contains a figure environment, an image-inclusion command or drawing code, so no image is delivered and the package declares no asset dependency.  Editorial formatting only, in addition: an AMS \texttt{amsart} preamble that restates the article's own declarations and macros used by the retained lines, namely the environments \texttt{lemma} on separate counters, together with the standard packages those lines need.  The retained environments print in this document as Lemma 1 in order of appearance, whereas the article prints the same items with the numbering of its own full document.  The complete native source of the article as distributed in the arXiv e-print for this exact version is bundled unchanged as \texttt{sources/194915/source0.tex}.
\begin{lemma}\label[lemma]{lemma:affine-series-ratio}
Let \(\alpha,\beta>0\), and let
\[
\mathbb D:=\{x\in\mathbb C:\ |x|<1\},\qquad \mathbb D^\ast:=\mathbb D\setminus(-1,0].
\]
For \(m\in\mathbb N\) and \(x\in\mathbb D^\ast\), define
\[
T_m^{(\alpha,\beta)}(x):=\sum_{n\ge0}(\alpha n+\beta)^m x^n.
\]
Then
\[
\frac1m\,\frac{T_{m+1}^{(\alpha,\beta)}(x)}{T_m^{(\alpha,\beta)}(x)}
\longrightarrow
\frac{\alpha}{-\log x},
\qquad m\to\infty,
\]
locally uniformly on \(\mathbb D^\ast\), where \(\log\) denotes the principal branch.
\end{lemma}

\begin{proof}
Fix a compact set \(K\subset\mathbb D^\ast\). We prove the convergence uniformly on \(K\).

For \(m\in\mathbb N\),
\[
T_m^{(\alpha,\beta)}(x)=\sum_{n\ge0}(\alpha n+\beta)^m x^n.
\]
By Cauchy's integral formula for derivatives, for any \(r>0\) and \(n\ge0\),
\[
(\alpha n+\beta)^m
=
\left.\frac{d^m}{dz^m}e^{(\alpha n+\beta)z}\right|_{z=0}
=
\frac{m!}{2\pi i}\oint_{|z|=r}\frac{e^{(\alpha n+\beta)z}}{z^{m+1}}\,dz.
\]
Hence
\[
T_m^{(\alpha,\beta)}(x)
=
\frac{m!}{2\pi i}\oint_{|z|=r}\frac{1}{z^{m+1}}
\sum_{n\ge0}x^n e^{(\alpha n+\beta)z}\,dz.
\]

Since \(K\subset\mathbb D\) is compact, there exists \(M<1\) such that \(|x|\le M\) for all \(x\in K\).
Choose \(r>0\) such that
\[
Me^{\alpha r}<1.
\]
Then, for \(|z|=r\) and \(x\in K\),
\[
|x e^{\alpha z}|
\le |x|e^{\alpha|z|}
\le Me^{\alpha r}
<1,
\]
so the series \(\sum_{n\ge0}x^n e^{(\alpha n+\beta)z}\) converges absolutely and uniformly on
\(K\times\{|z|=r\}\). Therefore we may interchange sum and integral, obtaining
\[
\sum_{n\ge0}x^n e^{(\alpha n+\beta)z}
=
e^{\beta z}\sum_{n\ge0}(x e^{\alpha z})^n
=
\frac{e^{\beta z}}{1-x e^{\alpha z}}.
\]
Thus
\[
T_m^{(\alpha,\beta)}(x)
=
\frac{m!}{2\pi i}\oint_{|z|=r}\frac{g_x(z)}{z^{m+1}}\,dz,
\qquad
g_x(z):=\frac{e^{\beta z}}{1-x e^{\alpha z}}.
\]

The poles of \(g_x\) are the solutions of \(1-x e^{\alpha z}=0\), namely
\[
z_k(x)=\frac{-\log x+2\pi i k}{\alpha},
\qquad k\in\mathbb Z.
\]
Write
\[
\log x=\log|x|+i\Arg(x),
\qquad -\pi<\Arg(x)<\pi,
\]
and set
\[
z_0(x):=\frac{-\log x}{\alpha}.
\]
Since \(x\in\mathbb D^\ast\), we have \(-\log x\neq0\). Moreover, because \(K\subset\mathbb D^\ast\) is compact,
there exists \(b<\pi\) such that
\[
|\Arg(x)|\le b
\qquad (x\in K).
\]
Hence, for \(k\neq0\),
\[
|2\pi k-\Arg(x)|\ge 2\pi-b>b\ge |\Arg(x)|.
\]
It follows that
\[
|z_k(x)|^2
=
\frac{(-\log|x|)^2+(2\pi k-\Arg(x))^2}{\alpha^2}
>
\frac{(-\log|x|)^2+\Arg(x)^2}{\alpha^2}
=
|z_0(x)|^2.
\]
Thus \(z_0(x)\) is the unique pole of \(g_x\) of minimal modulus, uniformly for \(x\in K\).

Set
\[
d_0:=\sup_{x\in K}|z_0(x)|,
\qquad
d_1:=\inf_{\substack{x\in K\\ k\neq0}}|z_k(x)|.
\]
Then \(0<d_0<d_1\). Choose \(R>0\) such that
\[
d_0<R<d_1.
\]
For each \(x\in K\), the function \(g_x\) is meromorphic on and inside \(|z|=R\), with exactly one pole
in the closed annulus \(r\le |z|\le R\), namely \(z_0(x)\). By the residue theorem,
\[
\oint_{|z|=r}\frac{g_x(z)}{z^{m+1}}\,dz
=
\oint_{|z|=R}\frac{g_x(z)}{z^{m+1}}\,dz
-
2\pi i\,\operatorname{Res}\!\left(\frac{g_x(z)}{z^{m+1}},z_0(x)\right).
\]

The pole at \(z_0(x)\) is simple, and its residue is
\[
\operatorname{Res}(g_x,z_0(x))
=
\frac{e^{\beta z_0(x)}}{(1-x e^{\alpha z})'|_{z=z_0(x)}}
=
\frac{e^{\beta z_0(x)}}{-\alpha x e^{\alpha z_0(x)}}
=
-\frac{e^{\beta z_0(x)}}{\alpha},
\]
since \(x e^{\alpha z_0(x)}=1\). Therefore,
\[
\operatorname{Res}\!\left(\frac{g_x(z)}{z^{m+1}},z_0(x)\right)
=
-\frac{e^{\beta z_0(x)}}{\alpha\,z_0(x)^{m+1}},
\]
and hence
\[
T_m^{(\alpha,\beta)}(x)
=
\frac{m!}{2\pi i}\oint_{|z|=R}\frac{g_x(z)}{z^{m+1}}\,dz
+
\frac{m!e^{\beta z_0(x)}}{\alpha\,z_0(x)^{m+1}}.
\]

Now \(1-xe^{\alpha z}\neq0\) on \(K\times\{|z|=R\}\), and this set is compact. Hence
\[
|g_x(z)|\le C
\qquad (x\in K,\ |z|=R)
\]
for some constant \(C>0\). Thus
\[
\left|
\frac{m!}{2\pi i}\oint_{|z|=R}\frac{g_x(z)}{z^{m+1}}\,dz
\right|
\le
\frac{m!}{2\pi}(2\pi R)\frac{C}{R^{m+1}}
=
m!CR^{-m},
\]
uniformly for \(x\in K\).

Since \(K\) is compact and \(0\notin K\), there exists \(c>0\) such that
\[
|z_0(x)|\ge c
\qquad (x\in K).
\]
Therefore
\[
T_m^{(\alpha,\beta)}(x)
=
\frac{m!e^{\beta z_0(x)}}{\alpha\,z_0(x)^{m+1}}
\left(
1+O\!\left(\left(\frac{d_0}{R}\right)^m\right)
\right),
\qquad m\to\infty,
\]
uniformly for \(x\in K\). Applying the same estimate with \(m+1\) in place of \(m\), we get
\[
T_{m+1}^{(\alpha,\beta)}(x)
=
\frac{(m+1)!e^{\beta z_0(x)}}{\alpha\,z_0(x)^{m+2}}
\left(
1+O\!\left(\left(\frac{d_0}{R}\right)^{m+1}\right)
\right),
\]
uniformly on \(K\). Since \(d_0/R<1\), it follows that
\[
\frac{T_{m+1}^{(\alpha,\beta)}(x)}{T_m^{(\alpha,\beta)}(x)}
=
\frac{m+1}{z_0(x)}\bigl(1+o(1)\bigr),
\]
uniformly for \(x\in K\). Dividing by \(m\), we obtain
\[
\frac1m\,\frac{T_{m+1}^{(\alpha,\beta)}(x)}{T_m^{(\alpha,\beta)}(x)}
\longrightarrow
\frac1{z_0(x)}
=
\frac{\alpha}{-\log x},
\]
uniformly on \(K\). Since \(K\subset\mathbb D^\ast\) was arbitrary, the convergence is locally uniform on \(\mathbb D^\ast\).
This proves the claim.
\end{proof}

\end{document}
